%=============================================================
% CHAPTER 1: INTRODUCTION
%=============================================================

Named Entity Recognition (NER) is a fundamental task in Natural Language Processing (NLP) that involves identifying and classifying named entities in text, such as persons, organizations, and locations \cite{mohit2014ner}. By extracting structured information from unstructured texts, NER facilitates the identification of relevant information and supports a wide range of NLP applications, including information extraction, document classification, and question answering. It is particularly useful for processing large volumes of textual data, allowing relevant entities to be identified more efficiently.

Recent advances in NLP have been driven by Transformer-based models, which have demonstrated strong performance across a wide range of language understanding tasks \cite{vaswani2017attention}. Pretrained language models can be fine-tuned on annotated datasets to adapt their representations to specific tasks, including NER. However, model performance may vary depending on the language and the characteristics of the text being processed.

This project investigates the performance of two pretrained Transformer models on Italian NER: UmBERTo, a model pretrained specifically for Italian, and XLM-RoBERTa, a multilingual model trained on text from multiple languages. The experiments use the NERMuD 2023 dataset, which covers three distinct domains: news, fiction, and political speeches. This dataset provides an opportunity to evaluate the models across different types of Italian text.

The main goal is to compare the performance of these two models under comparable experimental conditions and investigate whether language-specific pretraining provides an advantage over multilingual pretraining for Italian NER. Both models are fine-tuned on the same training data and evaluated using the same evaluation procedure.

The remainder of this report is organized as follows. Chapter 1 presents relevant background and related work. Chapter 2 describes the dataset and experimental methodology. Chapter 3 presents the experimental results followed by a discussion. Finally, Chapter 4 summarizes the main conclusions.

%-------------------------------------------------------------
\section{Background and Related work}
\label{sec:background}
%-------------------------------------------------------------

\subsection{Named Entity Recognition}
\label{subsec:ner}

Named Entity Recognition (NER) is an NLP task that consists of identifying and classifying named entities in text into predefined categories \cite{mohit2014ner}. Common entity categories include persons, organizations, and locations.

For example, consider the following sentence: ``The operation for Pope Leo will take place at the Agostino Gemelli Polyclinic Hospital in Rome.'' A NER system could identify \textit{Pope Leo} as a person, \textit{Agostino Gemelli Polyclinic Hospital} as an organization, and \textit{Rome} as a location. This process transforms unstructured text into structured information that can be used by other NLP applications.

NER is useful in a variety of applications, including information extraction, question answering, machine translation, and document analysis. By automatically identifying relevant entities, NER systems can facilitate the retrieval and organization of information from large collections of text.

However, identifying named entities can be challenging because the same word may refer to different entities depending on its context, and entities may consist of multiple words. These challenges make NER an important task in NLP research.

\subsection{Transformer-based Language Models}
\label{subsec:transformer}

Before the introduction of Transformers, recurrent neural networks (RNNs), particularly Bidirectional Long Short-Term Memory networks (BLSTMs) combined with Conditional Random Fields (CRFs), were widely used for sequence labeling tasks such as Named Entity Recognition. The introduction of the Transformer architecture marked a significant development in NLP, enabling models to capture contextual relationships between words through a mechanism known as self-attention \cite{vaswani2017attention}.

Unlike recurrent architectures, which process sequences step by step, Transformers can process tokens in parallel during training and model relationships between tokens regardless of their distance in a sequence. This allows them to learn contextual representations that are useful for a wide range of language understanding tasks.

The development of pretrained Transformer-based language models, such as BERT, further advanced NLP by enabling models to learn general linguistic representations from large text collections. These models can then be fine-tuned on task-specific annotated datasets, adapting their learned representations to tasks such as NER \cite{devlin2019bert}.

This approach is particularly relevant to the present study, which investigates whether an Italian-specific pretrained model, UmBERTo, achieves better NER performance than the multilingual XLM-RoBERTa model. Both models are based on Transformer architectures and can be fine-tuned for token-level classification.


\subsection{umBERTo and XLM-RoBERTa}
\label{subsec:berts}

UmBERTo and XLM-RoBERTa are pretrained language models based on the Transformer architecture, but they differ in their pretraining data and language coverage.

UmBERTo is a RoBERTa-based language model specifically developed for Italian. The model used in this project, \texttt{umberto-commoncrawl-cased-v1}, was pretrained on a large Italian text corpus derived from OSCAR, which is based on Common Crawl data. It uses SentencePiece tokenization and Whole Word Masking, an approach in which entire words are masked during pretraining rather than individual subword tokens \cite{musixmatch2020umberto}.

In contrast, XLM-RoBERTa (XLM-R) is a multilingual language model pretrained on text from 100 languages using a large collection of filtered Common Crawl data. It follows a masked language modelling objective, learning contextual representations by predicting masked tokens. Its multilingual pretraining allows it to represent text in different languages within a shared model \cite{conneau2020xlmr}.

These differences make the two models suitable for comparison in Italian Named Entity Recognition. UmBERTo benefits from pretraining focused specifically on Italian, whereas XLM-RoBERTa learns representations from a multilingual corpus. Comparing their performance on the same dataset can help investigate whether language-specific pretraining provides an advantage for this task.

\subsection{NERMuD 2023}
\label{subsec:nermud}

NERMuD (Named-Entities Recognition on Multi-Domain Documents) was a shared task organized as part of the EVALITA 2023 evaluation campaign for Italian Natural Language Processing. Its objective was to evaluate named entity recognition systems on documents from different domains, focusing on three entity categories: persons, organizations, and locations \cite{palmeroaprosio2023nermud}.

The task covered three domains: news, fiction, and political speeches. It included two subtasks: Domain-Agnostic Classification (DAC), which required a single model to process documents from all three domains, and Domain-Specific Classification (DSC), which allowed a separate model to be developed for each domain \cite{palmeroaprosio2023nermud}.

In this project, the NERMuD 2023 dataset is used to fine-tune and evaluate UmBERTo and XLM-RoBERTa. By considering documents from different domains, the experiments investigate how the two models perform on Italian text beyond a single type of document. The comparison uses the same dataset splits and evaluation procedure for both models to support a fair assessment of their performance.